\documentclass{amsart}
% For dealing with references we use the comment environment
\usepackage{verbatim}
\newenvironment{reference}{\comment}{\endcomment}
%\newenvironment{reference}{}{}
\newenvironment{slogan}{\comment}{\endcomment}
\newenvironment{history}{\comment}{\endcomment}

% For commutative diagrams we use Xy-pic
\usepackage[all]{xy}

% We use 2cell for 2-commutative diagrams.
\xyoption{2cell}
\UseAllTwocells

% We use multicol for the list of chapters between chapters
\usepackage{multicol}

% This is generally recommended for better output
\usepackage{lmodern}
\usepackage[T1]{fontenc}

% For cross-file-references

% Package for hypertext links:
\usepackage{hyperref}

% For any local file, say "hello.tex" you want to link to please

% Theorem environments.
%
\theoremstyle{plain}
\newtheorem{theorem}[subsection]{Theorem}
\newtheorem{proposition}[subsection]{Proposition}
\newtheorem{lemma}[subsection]{Lemma}

\theoremstyle{definition}
\newtheorem{definition}[subsection]{Definition}
\newtheorem{example}[subsection]{Example}
\newtheorem{exercise}[subsection]{Exercise}
\newtheorem{situation}[subsection]{Situation}

\theoremstyle{remark}
\newtheorem{remark}[subsection]{Remark}
\newtheorem{remarks}[subsection]{Remarks}

\numberwithin{equation}{subsection}

% Macros
%
\def\lim{\mathop{\mathrm{lim}}\nolimits}
\def\colim{\mathop{\mathrm{colim}}\nolimits}
\def\Spec{\mathop{\mathrm{Spec}}}
\def\Hom{\mathop{\mathrm{Hom}}\nolimits}
\def\Ext{\mathop{\mathrm{Ext}}\nolimits}
\def\SheafHom{\mathop{\mathcal{H}\!\mathit{om}}\nolimits}
\def\SheafExt{\mathop{\mathcal{E}\!\mathit{xt}}\nolimits}
\def\Sch{\mathit{Sch}}
\def\Mor{\mathop{\mathrm{Mor}}\nolimits}
\def\Ob{\mathop{\mathrm{Ob}}\nolimits}
\def\Sh{\mathop{\mathit{Sh}}\nolimits}
\def\NL{\mathop{N\!L}\nolimits}
\def\CH{\mathop{\mathrm{CH}}\nolimits}
\def\proetale{{pro\text{-}\acute{e}tale}}
\def\etale{{\acute{e}tale}}
\def\QCoh{\mathit{QCoh}}
\def\Ker{\mathop{\mathrm{Ker}}}
\def\Im{\mathop{\mathrm{Im}}}
\def\Coker{\mathop{\mathrm{Coker}}}
\def\Coim{\mathop{\mathrm{Coim}}}

% Boxtimes
%
\DeclareMathSymbol{\boxtimes}{\mathbin}{AMSa}{"02}

%
% Macros for moduli stacks/spaces
%
\def\QCohstack{\mathcal{QC}\!\mathit{oh}}
\def\Cohstack{\mathcal{C}\!\mathit{oh}}
\def\Spacesstack{\mathcal{S}\!\mathit{paces}}
\def\Quotfunctor{\mathrm{Quot}}
\def\Hilbfunctor{\mathrm{Hilb}}
\def\Curvesstack{\mathcal{C}\!\mathit{urves}}
\def\Polarizedstack{\mathcal{P}\!\mathit{olarized}}
\def\Complexesstack{\mathcal{C}\!\mathit{omplexes}}
% \Pic is the operator that assigns to X its picard group, usage \Pic(X)
% \Picardstack_{X/B} denotes the Picard stack of X over B
% \Picardfunctor_{X/B} denotes the Picard functor of X over B
\def\Pic{\mathop{\mathrm{Pic}}\nolimits}
\def\Picardstack{\mathcal{P}\!\mathit{ic}}
\def\Picardfunctor{\mathrm{Pic}}
\def\Deformationcategory{\mathcal{D}\!\mathit{ef}}

\setlength{\emergencystretch}{3em}
\begin{document}
\title[Stacks Project, Tag 0AVB]{337 --- Reflexive modules and codimension-one intersections}
\maketitle
\noindent Copyright (C) 2005--2025 Johan de Jong. Stacks Project authors. Source: \href{https://stacks.math.columbia.edu/tag/0AVB}{Tag 0AVB}.

\noindent Editorial difficulty note (not author text). Difficulty H2: The key direction expresses the codimension-one intersection inside the generic fibre as a kernel with a direct-sum target. Localizing this description and using depth at higher-codimension primes proves equality with the original module, after which DVR reflexivity identifies the double dual..

\noindent Editorial provenance note (not author text). History: excerpt from revision \texttt{a0444\allowbreak 6e57e\allowbreak c1fbc\allowbreak 252a8\allowbreak 71afc\allowbreak ec775\allowbreak 2fb28\allowbreak 07b14}, more-algebra.tex, lines 5577--5633. Reference presentation was adapted; original-author context and core support blocks were retained or added. The mathematical wording is unchanged. Licensed under GNU FDL 1.2 or later; no invariant sections or cover texts. See ../licenses/COPYING. Context blocks reproduce author definitions and supporting results; they are not additional counted problems.

\begin{definition}
\label{definition-reflexive}
Let $R$ be a domain. We say an $R$-module $M$ is {\it reflexive} if
the natural map
$$
j : M \longrightarrow \Hom_R(\Hom_R(M, R), R)
$$
which sends $m \in M$ to the map sending $\varphi \in \Hom_R(M, R)$
to $\varphi(m) \in R$ is an isomorphism.
\end{definition}

\begin{lemma}
\label{lemma-reflexive-over-normal}
Let $R$ be a Noetherian normal domain with fraction field $K$.
Let $M$ be a finite $R$-module. The following are equivalent
\begin{enumerate}
\item $M$ is reflexive,
\item $M$ is torsion free and has property $(S_2)$,
\item $M$ is torsion free and
$M = \bigcap_{\text{height}(\mathfrak p) = 1} M_{\mathfrak p}$
where the intersection happens in $M_K = M \otimes_R K$.
\end{enumerate}
\end{lemma}

\begin{proof}
By Algebra, Lemma \href{https://stacks.math.columbia.edu/tag/031S}{lemma criterion normal (Tag 031S)}
we see that $R$ satisfies $(R_1)$ and $(S_2)$.

\medskip\noindent
Assume (1). Then $M$ is torsion free by Lemma \href{https://stacks.math.columbia.edu/tag/0AV0}{lemma reflexive torsion free (Tag 0AV0)}
and satisfies $(S_2)$ by Lemma \href{https://stacks.math.columbia.edu/tag/0EBA}{lemma reflexive S2 (Tag 0EBA)}.
Thus (2) holds.

\medskip\noindent
Assume (2). By definition
$M' = \bigcap_{\text{height}(\mathfrak p) = 1} M_{\mathfrak p}$
is the kernel of the map
$$
M_K
\longrightarrow
\bigoplus\nolimits_{\text{height}(\mathfrak p) = 1} M_K/M_\mathfrak p
\subset
\prod\nolimits_{\text{height}(\mathfrak p) = 1} M_K/M_\mathfrak p
$$
Observe that our map indeed factors through the direct sum as indicated
since given
$a/b \in K$ there are at most finitely many height $1$ primes $\mathfrak p$
with $b \in \mathfrak p$. Let $\mathfrak p_0$ be a prime of height $1$.
Then $(M_K/M_\mathfrak p)_{\mathfrak p_0} = 0$ unless
$\mathfrak p = \mathfrak p_0$ in which case we get
$(M_K/M_\mathfrak p)_{\mathfrak p_0} = M_K/M_{\mathfrak p_0}$.
Thus by exactness of localization and the fact that localization
commutes with direct sums, we see that
$M'_{\mathfrak p_0} = M_{\mathfrak p_0}$.
Since $M$ has depth $\geq 2$ at primes of height $> 1$,
we see that $M \to M'$ is an isomorphism by
Lemma \href{https://stacks.math.columbia.edu/tag/0AV9}{lemma isom depth 2 torsion free (Tag 0AV9)}. Hence (3) holds.

\medskip\noindent
Assume (3). Let $\mathfrak p$ be a prime of height $1$.
Then $R_\mathfrak p$ is a discrete valuation ring by $(R_1)$.
By Lemma \href{https://stacks.math.columbia.edu/tag/0AUW}{lemma dedekind torsion free flat (Tag 0AUW)}
we see that $M_\mathfrak p$ is finite free, in particular
reflexive. Hence the map $M \to M^{**}$ induces an isomorphism at all
the primes $\mathfrak p$ of height $1$. Thus the condition
$M = \bigcap_{\text{height}(\mathfrak p) = 1} M_{\mathfrak p}$
implies that $M = M^{**}$ and (1) holds.
\end{proof}
\end{document}
